\chapter{Stress Equilibrium Equation}

\section*{Introduction}

Imagine a 3D grid of springs connecting every atom to its neighbors. When you pull on one end, the tension propagates through the spring network. The stress equilibrium equation says that at each node, the net force from all connected springs equals zero (in statics) or equals mass times acceleration (in dynamics). Engineers use this to predict where bridges will crack, where airplane wings will fatigue, and where earthquake forces will concentrate in buildings.


\[
\frac{\partial \sigma_{ij}}{\partial x_j} + f_i = \rho \frac{\partial^2 u_i}{\partial t^2}
\]

\section*{Fundamental Equations Used}

The derivation starts from force balance on an infinitesimal cube using the components of the stress tensor.

\section*{Assumptions}

\textbf{Assumptions:} The material obeys Newton's second law. The stress tensor $\sigma_{ij}$ is defined at every point. Body forces (gravity, electromagnetic) are included as $f_i$.


\textbf{Limitations:} Does not describe material failure (need fracture mechanics). Assumes continuous media---breaks down at atomic scales. Does not account for residual stresses from manufacturing.


\section*{Mathematical Derivation}

\textbf{Step 1: Consider an infinitesimal cube of material.}

Take a small cubic volume element of side lengths $dx$, $dy$, $dz$ at position $(x, y, z)$. The stress tensor $\sigma_{ij}$ describes the force per unit area on each face. On the face perpendicular to the $x$-axis at position $x$, the traction (force per unit area) is $\sigma_{ix}$; on the opposite face at $x + dx$, it is $\sigma_{ix} + (\partial\sigma_{ix}/\partial x)\,dx$.

\textbf{Step 2: Sum forces in the $x$-direction.}

The net force on the cube in the $x$-direction from the two faces perpendicular to $x$ is
\begin{equation}
\left(\sigma_{xx} + \frac{\partial\sigma_{xx}}{\partial x}dx\right)dy\,dz - \sigma_{xx}\,dy\,dz = \frac{\partial\sigma_{xx}}{\partial x}\,dx\,dy\,dz.
\end{equation}
Including contributions from the faces perpendicular to $y$ and $z$ (which carry $\sigma_{xy}$ and $\sigma_{xz}$ tractions):
\begin{equation}
\text{Net surface force in }x = \left(\frac{\partial\sigma_{xx}}{\partial x} + \frac{\partial\sigma_{xy}}{\partial y} + \frac{\partial\sigma_{xz}}{\partial z}\right)dx\,dy\,dz = \frac{\partial\sigma_{xj}}{\partial x_j}\,dV,
\end{equation}
using Einstein summation convention over $j$.

\textbf{Step 3: Add the body force.}

If a body force $f_i$ (force per unit volume, e.g., gravity $\rho g_i$) acts on the cube, the total force is
\begin{equation}
F_i = \left(\frac{\partial\sigma_{ij}}{\partial x_j} + f_i\right)dV.
\end{equation}

\textbf{Step 4: Apply Newton's second law.}

By Newton's second law, $F_i = (dm)\,a_i = \rho\,dV\,\partial^2 u_i/\partial t^2$, where $u_i$ is the displacement and $\rho$ is the mass density:
\begin{equation}
\frac{\partial\sigma_{ij}}{\partial x_j} + f_i = \rho\,\frac{\partial^2 u_i}{\partial t^2}.
\end{equation}

\textbf{Step 5: Write the final result.}

\begin{equation}
\boxed{\frac{\partial\sigma_{ij}}{\partial x_j} + f_i = \rho\,\frac{\partial^2 u_i}{\partial t^2}.}
\end{equation}
In index-free notation: $\nabla\cdot\boldsymbol{\sigma} + \mathbf{f} = \rho\,\ddot{\mathbf{u}}$. For statics ($\ddot{\mathbf{u}} = 0$), this reduces to $\nabla\cdot\boldsymbol{\sigma} + \mathbf{f} = \mathbf{0}$. The equation represents local force balance at every point in the material.

\section*{Characteristics of the Equation}

The equation is a set of 3 coupled first-order PDEs (in terms of stress) or second-order PDEs (in terms of displacement via compatibility). It applies in Cartesian, cylindrical, and spherical coordinates. For statics, it reduces to $\partial \sigma_{ij}/\partial x_j + f_i = 0$.
\section*{Solution Methods}

For simple geometries: direct integration with boundary conditions. For complex geometries: finite element analysis (FEA), finite difference methods, boundary element methods. Airy stress function approach reduces 2D problems to a single biharmonic equation.
\section*{Verifications}

Strain gauges measure local deformation, which can be converted to stress. Photoelasticity visualizes stress fields in transparent models. FEA results are validated against analytical solutions for simple cases (thick-walled cylinder, half-space loading).
\section*{Applications}

Structural engineering (stress analysis of buildings, bridges, aircraft), geomechanics (stress in rock formations, borehole stability), biomechanics (stress in bones, blood vessels), manufacturing (stamping, forging simulations).
\section*{What Richard Feynman Would Have Asked}

\subsection*{1. The Plain-English Translation}
\textit{``Don't just give me the tensor equation. What is this really saying about the material?''}

The Core Meaning: Every tiny cube of material is in a tug-of-war with its neighbors. The stress equilibrium equation says that if the cube is not accelerating, the total pull from all six faces plus any body forces (gravity) must exactly cancel. It is Newton's second law applied to an infinitesimal element, and it works for any material---steel, rubber, concrete, or blood.

The Metaphor: Think of a chain of people holding ropes in every direction. If nobody is moving, then the total pull on each person from all their ropes must be zero. The stress equilibrium equation says the same thing for every point inside a material: all the ``pulls'' (stresses) from neighboring points must balance.

\subsection*{2. The Localized Mechanics}
\textit{``Zoom into a single point inside the material. What forces are acting there?''}

He would press you on what the stress tensor actually represents. At any point, there are three planes (perpendicular to $x$, $y$, $z$), and on each plane there are three force components---giving nine stress components (reduced to six by symmetry). The equilibrium equation says that the gradient of these stresses must balance the body force. If $\partial \sigma_{xx}/\partial x$ is positive, more tension is pulling to the right, and something must balance it.

The Smoothness Check: He would ask whether the stress field must be continuous. In general, yes---but at interfaces between different materials, stress can be discontinuous (though the traction vector must be continuous). At crack tips, stresses become singular---the equilibrium equation still holds, but the solution becomes infinite, signaling the need for fracture mechanics.

\subsection*{3. Testing the Extreme Limits}
\textit{``What happens when we push the material to extreme conditions?''}

The Hydrostatic Limit: Under pure hydrostatic pressure ($\sigma_{ij} = -P\delta_{ij}$), the equilibrium equation reduces to $\nabla P = \rho g$---the familiar equation for pressure in a fluid. This is why submarines must withstand enormous pressure: every point inside the hull is being squeezed equally from all directions.

The Crack-Tip Limit: Near a sharp crack, stresses scale as $1/\sqrt{r}$ where $r$ is the distance from the tip. The equilibrium equation is satisfied everywhere except at the tip itself, where it breaks down. This singularity drives crack propagation and is the foundation of linear elastic fracture mechanics.

The Atomic Scale: At scales below $\sim 10$ nm, the concept of ``stress at a point'' breaks down because there are too few atoms to define a meaningful average. Molecular dynamics simulations replace the continuum equation with direct atom-by-atom force calculations.

\subsection*{4. Dimensionality and Geometry}
\textit{``How does the geometry of the structure affect the stress distribution?''}

He would point out that the same equilibrium equation applies to a skyscraper, a Dam, and a contact lens---the difference is entirely in the boundary conditions and geometry. In 2D (plane stress or plane strain), the equilibrium equation has only 2 components. In 3D, it has 3. The mathematical structure is identical; the complexity comes from the geometry of the boundaries.

\subsection*{5. Cross-Disciplinary Connection}
\textit{``Where else does this exact same force-balance equation appear?''}

The same equation structure appears in: momentum balance in fluid mechanics (Navier-Stokes), where the stress tensor includes viscous terms; in electromagnetism (force balance on charged fluids, magnetohydrodynamics); in biological tissues (tumor growth mechanics, where ``stress'' includes cellular contractile forces); and in economics (``stress'' in supply chains, where the equilibrium equation becomes a balance of flows).
\section*{FAQ}

\textbf{Q: What is the difference between stress equilibrium and strain compatibility?}\\
\textbf{A:} Stress equilibrium relates forces to stresses (statics). Strain compatibility ensures that the deformation is geometrically consistent (no gaps or overlaps). Together, they determine the complete stress-strain state.

\textbf{Q: Can the stress equilibrium equation predict when a material will break?}\\
\textbf{A:} No. It describes the state of stress in an intact material. Failure criteria (von Mises, Tresca, Mohr-Coulomb) are needed to predict when the stress exceeds the material's strength.

\textbf{Q: Why is the stress tensor symmetric?}\\
\textbf{A:} From conservation of angular momentum. If the stress tensor were asymmetric, there would be a net torque on an infinitesimal element, causing angular acceleration---which contradicts the equilibrium of moments in the absence of body couples.
\section*{Important Bibliography}

\begin{enumerate}
\item Timoshenko, S.P. and Goodier, J.N., \textit{Theory of Elasticity}, 3rd ed. (McGraw-Hill, 1970), Chapter 2.
\item Fung, Y.C., \textit{A First Course in Continuum Mechanics}, 3rd ed. (Prentice Hall, 1994), Chapter 5.
\item Sadd, M.H., \textit{Elasticity: Theory, Applications, and Numerics}, 3rd ed. (Academic Press, 2014), Chapter 2.
\item Malvern, L.E., \textit{Introduction to the Mechanics of a Continuous Medium} (Prentice Hall, 1969), Chapter 5.
\end{enumerate}

